% !PRINTED 582
\weberpage{582}

Now for variable $\varkappa$ and $v$ we have [\S\,45, (5)]:
\weq{8}{\mathrm{dn}\,(v,\varkappa) = \dfrac{\mathrm{dn}\,(iv,\varkappa')}{\mathrm{cn}\,(iv,\varkappa')}.}
If we here set:
\[
v = \Omega_{h,h'} = \dfrac{2hK + 2h'iK'}{m},
\]
\[
iv = \dfrac{-2h'K' + 2hiK}{m} = \Omega'_{-h',h},
\]
then the $\Omega'$ have the same meaning for $\varkappa'$ as the $\Omega$ have for $\varkappa$, i.\,e. they correspond to the interchange of $(\omega, -1/\omega)$, and it follows from (7) and (8) that also
\[
\dfrac{\mathrm{dn}\,\Omega_{h,h'}}{\sqrt{\varkappa'}} \text{ is a unit.}
\]
Thus also
\weq{9}{\begin{aligned}
\dfrac{1}{\sqrt{\varkappa'}}\,\mathrm{dn}\,\Omega_{h,h'}
&= \dfrac{\vartheta_{00}\left(\dfrac{h+h'\omega}{m}\right)}{\vartheta_{01}\left(\dfrac{h+h'\omega}{m}\right)}, \\
\sqrt{\dfrac{\varkappa}{\varkappa'}}\,\mathrm{cn}\,\Omega_{h,h'}
&= \dfrac{\vartheta_{10}\left(\dfrac{h+h'\omega}{m}\right)}{\vartheta_{01}\left(\dfrac{h+h'\omega}{m}\right)}\ \text{are algebraic units.}
\end{aligned}}
From formula \S\,155, (20) there follows for $v=0$:
\weq{10}{\pm m = \overset{h,h'}{\varPi}\ \sqrt{\varkappa}\ \mathrm{sn}\,\Omega_{h,h'} = \varPi\, \varkappa_{h,h'}.}
1. \gesperrt{The integers $\varkappa_{\mu,\mu'}$ are therefore divisors of the number $m$.}

We now set, when $m$ and $n$ are two odd numbers:
\[
\Omega_{h,h'} = \dfrac{2hK + 2h'iK'}{m},
\]
\[
H_{l,v'} = \dfrac{2lK + 2l'iK'}{n}, % [source sic: the printed subscript is v' although the numerator uses l']
\]
and substitute in formula \S\,155, (20):
\[
v = H_{l,v'}.
\]
This gives:
\weq{11}{\dfrac{\sqrt{\varkappa}\,\mathrm{sn}\,m\,H_{l,v'}}{\sqrt{\varkappa}\,\mathrm{sn}\,H_{l,v'}} = \overset{h,h'}{\varPi}\ \sqrt{\varkappa}\ \mathrm{sn}\,(H_{l,v'} + \Omega_{h,h'}),
}
%CONT>>
